\documentclass[a4paper,11pt]{extarticle}

\usepackage{cursoice}

% ----------------------------------------------
% custom title
% ----------------------------------------------
\makeatletter
\renewcommand*{\maketitle}{%
\noindent
\begin{minipage}{0.2\textwidth}
	\includegraphics[width=\textwidth]{img/logo_dte.png}
\end{minipage}
\hspace{1em}
\begin{minipage}{0.56\textwidth}
	\begin{tikzpicture}
	\node[rectangle,rounded corners=6pt,inner sep=10pt,fill=blue!50!black,text width= 0.95\textwidth] {\color{white} {\Huge \@title} \\ {\Large}};
	\end{tikzpicture}
\end{minipage}
\hspace{1em}
\begin{minipage}{0.1\textwidth}
	\includegraphics[width=\textwidth]{img/logo_teleco_2.png}
\end{minipage}
\hfill
\bigskip\bigskip
}%
\makeatother
% ------------------------------------------------

% ------------------------------------------------
% custom section
% ------------------------------------------------
\usepackage[explicit]{titlesec}
\newcommand*\sectionlabel{}
\titleformat{\section}
  {\gdef\sectionlabel{}
   \normalfont\sffamily\Large\bfseries\scshape}
  {\gdef\sectionlabel{\thesection\ }}{0pt}
  {
\noindent
\begin{tikzpicture}
\node[rectangle,rounded corners=3pt,inner sep=4pt,fill=blue!50!black,text width= 0.95\columnwidth] {\color{white}\sectionlabel#1};
\end{tikzpicture}
  }
\titlespacing*{\section}{0pt}{15pt}{2pt}
% -----------------------------------------------

% -----------------------------------------------
% custom footer
% -----------------------------------------------
\usepackage{fancyhdr}
\makeatletter
\pagestyle{fancy}
\fancyhead{}
\fancyfoot[C]{\footnotesize  \@date\ \ \@author}
\renewcommand{\headrulewidth}{0pt}
\renewcommand{\footrulewidth}{0pt}
\makeatother
% -----------------------------------------------


\title{Introducción a la \\ programación funcional}
\author{Santiago Higuera. Universidad Politécnica de Madrid}
\date{\today}


\newcommand{\raya}{
	\vspace{0.5em}
	{\centering \rule{10cm}{0.1mm} \par}
	\vspace{0.5em}
}

\begin{document}

\maketitle

{\centering \footnotesize \textbf{\color{red}(Versión de fecha  \today)} \par}


%\setcounter{tocdepth}{1}
%\vspace*{3em}
%\tableofcontents

%\vfill
\begin{Resumen}
Se desarrrolla a continuación un ejercicio que de manera progresiva irá introduciendo los conceptos aprendidos en la parte teórica de este curso.

\end{Resumen}

\section{Introducción}
Para desarrollar el aprendizaje de la codificación de programas en lenguaje Rust y los principios de programación funcional que se han desarrollado en la parte teórica del curso, se plantea un ejercicio que se irá resolviendo por pasos sucesivos.

Se supone que se está desarrollando el software de gestión de una editorial de libros que recibe libros de los autores y los añade a su catálogo.

Cada libro se simboliza por una estructura cuyos campos son el título, que es un String, el identificador del autor y el número de páginas, ambos representados con un entero del tipo \textit{u32}. El código correspondiente sería el siguiente:

\begin{Codigo2}

\end{Codigo2}

\section{Creación del proyecto}
Se deberá situar en el directorio de sus proyectos y crear un nuevo proyecto Rust. Hemos decidido que el proyecto tenga la estructura de una librería por lo que la orden a ejecutar desde la consola será la siguiente:

\vspace{0.7em}
\begin{Codigo2}
cargo new --lib editorial
\end{Codigo2}

Si abre el proyecto con el IDE que esté utilizando y abre el fichero \textit{src/lib.rs} verá un código similar al siguiente:

\vspace{0.7em}
\begin{Codigo2}
pub fn add(left: usize, right: usize) -> usize {
   left + right
}

#[cfg(test)]
mod tests {
   use super::*;
	
   #[test]
   fn it_works() {
      let result = add(2, 2);
      assert_eq!(result, 4);
   }
}
\end{Codigo2}

La librería recién creada consta de una función \textit{add()}, que realiza la suma de dos números enteros \textit{usize} y un test que permite probar la función. Para ejecutar el test debe teclear:

\vspace{0.7em}
\begin{Codigo2}
cargo test --lib
\end{Codigo2}

Obtendrá la salida correspondiente a la ejecución del test que, si todo ha ido bien, indicará que el test ha dado resultado correcto:

\vspace{0.7em}
\begin{Codigo2}
running 1 test
test tests::it_works ... ok

test result: ok. 1 passed; 0 failed; 0 ignored;
\end{Codigo2}

\section{Creación de los tipos de datos usados por el proyecto}
A continuación, dentro del fichero \textit{lib.rs}, vamos a crear las estructuras que se utilizarán como tipos de datos del proyecto. Añada el siguiente código al principio del fichero \textit{lib.rs}:

\vspace{0.7em}
\begin{Codigo2}
#[derive(Debug)]
struct Libro {
    titulo: String,
    id_autor: u32,
    num_paginas: u32
}
#[derive(Debug)]
struct Autor {
    id_autor: u32,
    nombre: String,
}
\end{Codigo2}

Recuerde que la etiqueta \textit{\#[derive(Debug)]} permite imprimir las estructuras utilizando el especificador de formato \textit{{:?}}.

Podemos crear algún test que nos permita probar las estructuras recién creadas. Teclee los siguientes test dentro del módulo \textit{tests} del fichero \textit{lib.rs}:

\vspace{0.7em}
\begin{Codigo2}
#[test]
fn test_libro() {
    let libro = Libro{titulo: String::from("Programación en Rust"),
        id_autor: 1, num_paginas: 228};
    assert_eq!(libro.titulo, "Programación en Rust");
    assert_eq!(libro.id_autor, 1);
    assert_eq!(libro.num_paginas, 228);
}
#[test]
fn test_autor() {
    let autor = Autor {id_autor: 1,
        nombre: String::from("Santiago Higuera")};
    assert_eq!(autor.nombre, "Santiago Higuera");
    assert_eq!(autor.id_autor, 1);
}    
\end{Codigo2}

Puede aprovechar ya para borrar el código de la función \textit{add()} y el del test \textit{it\_workds()}. Recuerde que la etiqueta \textit{\#[test]} es necesaria para indicar que determinada función es un test y se debe ejecutar al utilizar la orden \textit{cargo test}. Puede comprobar los test recién creados con la orden:

\vspace{0.7em}
\begin{Codigo2}
cargo test --lib
\end{Codigo2}

Si todo ha ido bien, recibirá en la consola un mensaje indicándolo.





\bibliography{Mibiblioteca}
\end{document}
